% sec1_7_part2.tex — Section 1.7 (part 2): Testing, R², residuals (pp. 65--70)

\subsection*{Testing the Homogeneity of the Cost Function}

Before proceeding to the estimation of the restricted model~\eqref{eq:1.7.6}, in order
to test the homogeneity restriction $\beta_3 + \beta_4 + \beta_5 = 1$, we will first
estimate the unrestricted model~\eqref{eq:1.7.4}.  If one uses the data available in
printed form in Nerlove's paper, the OLS estimate of the equation is:
\begin{equation}
\begin{aligned}
  \log(TC_i) &= \underset{(1.8)}{-3.5} + \underset{(0.017)}{0.72}\,\log(Q_i)
    + \underset{(0.29)}{0.44}\,\log(p_{i1}) \\
  &\quad - \underset{(0.34)}{0.22}\,\log(p_{i2})
    + \underset{(0.10)}{0.43}\,\log(p_{i3}) \\[6pt]
  R^2 &= 0.926,\quad \text{mean of dep.\ variable} = 1.72, \\
  SER &= 0.392,\quad SSR = 21.552,\quad n = 145.
\end{aligned}
\label{eq:1.7.7}
\end{equation}
Here, numbers in parentheses are the standard errors of the OLS coefficient estimates.
Since $\beta_2 = 1/r$, the estimate of the degree of returns to scale implied by the OLS
coefficient estimates is about $1.4$ $(= 1/0.72)$.  The OLS estimate of
$\beta_4 = \alpha_2/r$ has the wrong sign.  As noted by Nerlove, there are reasons to
believe that $p_{i2}$, the rental price of capital, is poorly measured.  This may explain
why $b_4$ is so imprecisely determined (i.e., the standard error is large relative to the
size of the coefficient estimate) that one cannot reject the hypothesis that $\beta_4 = 0$
with a $t$-ratio of $-0.65$ $(= -0.22/0.34)$.\footnote{The consequence of measurement
error is not just that the coefficient of the variable measured with error is poorly
determined; it could also contaminate the coefficient estimates for all other regressors.
The appropriate context to address this problem is the large sample theory for endogenous
regressors in Chapter~3.}

To test the homogeneity restriction $\text{H}_0\colon \beta_3 + \beta_4 + \beta_5 = 1$,
we could write the hypothesis in the form $\vec{R}\bm{\beta} = \vec{r}$ with
$\vec{R} = (0,\,0,\,1,\,1,\,1)$ and $\vec{r} = 1$ and use the formula~(1.4.9) to
calculate the $F$-ratio.  Alternatively, we can use the $F$-ratio formula~(1.4.11).  The
unrestricted model producing $SSR_U$ is~\eqref{eq:1.7.4} and the restricted model
producing $SSR_R$ is~\eqref{eq:1.7.6}, which superimposes the null hypothesis on the
unrestricted model.  The OLS estimate of~\eqref{eq:1.7.6} is
\begin{equation}
\begin{aligned}
  \log\!\left(\frac{TC_i}{p_{i3}}\right)
  &= \underset{(0.88)}{-4.7} + \underset{(0.017)}{0.72}\,\log(Q_i) \\
  &\quad + \underset{(0.20)}{0.59}\,\log(p_{i1}/p_{i3})
    - \underset{(0.19)}{0.007}\,\log(p_{i2}/p_{i3}) \\[6pt]
  R^2 &= 0.932,\quad \text{mean of dep.\ var.} = -1.48, \\
  SER &= 0.39,\quad SSR = 21.640,\quad n = 145.
\end{aligned}
\label{eq:1.7.8}
\end{equation}

The $F$ test of the homogeneity restriction proceeds as follows.

\textbf{Step~1:} Using~(1.4.11), the $F$-ratio can be calculated as
\[
  \frac{(21.640 - 21.552)/1}{21.552/(145 - 5)} = 0.57.
\]

\textbf{Step~2:} Find the critical value.  The number of restrictions (equations) in
the null hypothesis is 1, and $K$ (the number of coefficients in the unrestricted
model) is 5.  So the degrees of freedom are 1 and $140$ $(= 145 - 5)$.  From the table
of $F$ distributions, the critical value is about 3.9.

\textbf{Step~3:} Thus, we can easily accept the homogeneity restriction, a very
comforting conclusion for those who take microeconomics seriously (like us).

\subsection*{Detour: A Cautionary Note on $R^2$}

The $R^2$ of 0.926 is surprisingly high for cross-section estimates, but some of the
explanatory power of the regression comes from the scale effect that total costs increase
with firm size.  To gauge the contribution of the scale effect on the $R^2$, subtract
$\log(Q_i)$ from both sides of~\eqref{eq:1.7.4} to obtain an equivalent cost function:
\begin{equation}
\begin{aligned}
  \log\!\left(\frac{TC_i}{Q_i}\right)
  &= \beta_1 + (\beta_2 - 1)\log(Q_i) \\
  &\quad + \beta_3\log(p_{i1}) + \beta_4\log(p_{i2})
    + \beta_5\log(p_{i3}) + \varepsilon_i.
\end{aligned}
\tag{\ref{eq:1.7.4}$'$}
\label{eq:1.7.4p}
\end{equation}
Here, the dependent variable is the average cost rather than total costs.  Application of
the OLS to~\eqref{eq:1.7.4p} using the same data yields
\begin{equation}
\begin{aligned}
  \log\!\left(\frac{TC_i}{Q_i}\right)
  &= \underset{(1.8)}{-3.5} - \underset{(0.017)}{0.28}\,\log(Q_i) \\
  &\quad + \underset{(0.29)}{0.44}\,\log(p_{i1})
    - \underset{(0.34)}{0.22}\,\log(p_{i2})
    + \underset{(0.10)}{0.43}\,\log(p_{i3}) \\[6pt]
  R^2 &= 0.695,\quad \text{mean of dep.\ var.} = -4.83, \\
  SER &= 0.392,\quad SSR = 21.552,\quad n = 145.
\end{aligned}
\label{eq:1.7.9}
\end{equation}
As you no doubt have anticipated, the output coefficient is now $-0.28$ $(= 0.72 - 1)$
with the standard errors and the other coefficient estimates unchanged.  The $R^2$
changes only because the dependent variable is different.  It is nonsense to say that the
higher $R^2$ makes~\eqref{eq:1.7.4} preferable to~\eqref{eq:1.7.4p}, because the two
equations represent the same model.  The point is: when comparing equations on the basis
of the fit, the equations must share the same dependent variable.

\subsection*{Testing Constant Returns to Scale}

As an application of the $t$-test, consider testing whether returns to scale are constant
($r = 1$).  We take the maintained hypothesis to be the restricted model~\eqref{eq:1.7.6}.
Because $\beta_2$ (the log output coefficient) equals 1 if and only if $r = 1$, the null
hypothesis is that $\text{H}_0\colon \beta_2 = 1$.  The $t$-test of constant returns to
scale proceeds as follows.

\textbf{Step~1:} Calculate the $t$-ratio for the hypothesis.  From the estimation of the
restricted model, we have $b_2 = 0.72$ with a standard error of 0.017, so
\[
  t\text{-ratio} = \frac{0.72 - 1}{0.017} = -16.
\]
Because the maintained hypothesis here is the restricted model~\eqref{eq:1.7.6}, $K$
(the number of coefficients) $= 4$.

\textbf{Step~2:} Look for the critical value in the $t(141)$ distribution.  If the size
of the test is 5 percent, the critical value is 1.98.

\textbf{Step~3:} Since the absolute value of the $t$-ratio is far greater than the
critical value, we reject the hypothesis of constant returns to scale.

\subsection*{Importance of Plotting Residuals}

The regression has a problem that cannot be seen from the estimated coefficients and
their standard errors.  Figure~1.7 plots the residuals against $\log(Q_i)$.  Notice two
things from the plot.  First, as output increases, the residuals first tend to be
positive, then negative, and again positive.  This strongly suggests that the degree of
returns to scale ($r$) is not constant as assumed in the log-linear specification.
Second, the residuals are more widely scattered for lower outputs, which is a sign of a
failure of the homoskedasticity assumption that the error variance does not depend on the
regressors.  To deal with these problems, Nerlove divided the sample of 145 firms into
five groups of 29, ordered by output, and estimated the model~\eqref{eq:1.7.6}
separately for each group.  This amounts to allowing all the coefficients (including
$\beta_2 = 1/r$) \emph{and} the error variance to differ across the five groups
differing in size.  Nerlove finds that returns to scale diminish steadily, from a high of
well over 2 to a low of slightly below 1, over the output range of the data.  In the
empirical exercise of this chapter, the reader is asked to replicate this finding and do
some further analysis using \textbf{dummy variables} and the weighted least squares.

\begin{figure}[htbp]
\centering
\includegraphics[width=0.9\textwidth,trim=50 400 50 400,clip]{figures/fig_1_7.png}
\caption{Plot of Residuals against Log Output}
\label{fig:1.7}
\end{figure}

\subsection*{Subsequent Developments}

One strand of the subsequent literature is concerned about generalizing the Cobb-Douglas
technology while maintaining the assumption of cost minimization.  An obvious alternative
to Cobb-Douglas is the Constant Elasticity of Substitution (CES) production function,
but it has two problems.  First, the cost function implied by the CES production function
is highly nonlinear (which, though, could be overcome by the use of nonlinear least
squares to be covered in Chapter~7).  Second, the CES technology implies a constant
degree of returns to scale.  One of Nerlove's main findings is that the degree varies
with output.  Christensen and Greene (1976) are probably the first to estimate the
technology parameters allowing for variable degrees of returns to scale.  Using the
\textbf{translog cost function} introduced by Christensen, Jorgenson, and Lau (1973),
they find that the significant scale economies evident in the 1955 data were mostly
exhausted by 1970, with most firms operating at much higher output levels where the AC
curve is essentially flat.  Their work will be examined in detail in Chapter~4.

Another issue is whether regulated firms minimize costs.  The influential paper by
Averch and Johnson (1962) argues that the practice by regulators to guarantee utilities
a ``fair rate of return'' on their capital stock distorts the choice of input levels.
Consequently, unless the fair rate of return is used in the calculation of $p_{i2}$, the
true technology parameters cannot be estimated from the cost function.

\medskip\noindent
\textsc{Questions for Review}

\begin{enumerate}
\item (Review of duality theory) Consult your favorite microeconomic textbook to
  remember how to derive the Cobb-Douglas cost function from the Cobb-Douglas
  production function.

\item (Change of units) In Nerlove's data, output is measured in kilowatt hours.
  If output were measured in megawatt hours, how would the estimated restricted
  regression change?

\item (Recovering technology parameters from regression coefficients) Show that the
  technology parameters $(\mu, \alpha_1, \alpha_2, \alpha_3)$ can be determined uniquely
  from the first four equations in~\eqref{eq:1.7.5} and the definition
  $r \equiv \alpha_1 + \alpha_2 + \alpha_3$.  (Do not use the fifth equation
  $\beta_5 = \alpha_3/r$.)

\item (Recovering left-out coefficients from restricted OLS) Calculate the restricted
  OLS estimate of $\beta_5$ from~\eqref{eq:1.7.8}.  How do you calculate the standard
  error of $b_5$ from the printout of the restricted OLS?  \textbf{Hint:} Write
  $b_5 = a + \vec{c}'\vec{b}$ for suitably chosen $a$ and $\vec{c}$ where $\vec{b}$
  here is $(b_1, \ldots, b_4)'$.  So
  $\Var(b_5 \mid \vec{X}) = \vec{c}'\,\Var(\vec{b} \mid \vec{X})\,\vec{c}$.

\item If you take $p_{i2}$ instead of $p_{i3}$ and subtract $\log(p_{i2})$ from both
  sides of~\eqref{eq:1.7.4}, how does the restricted regression look?  Without actually
  estimating it on Nerlove's data, can you tell from the estimated restricted regression
  in the text what the restricted OLS estimate of $(\beta_1, \ldots, \beta_5)$ will be?
  Their standard errors?  The $SSR$?  What about the $R^2$?

\item Why is the $R^2$ of 0.926 from the unrestricted model~\eqref{eq:1.7.7}
  \emph{lower} than the $R^2$ of 0.932 from the restricted model~\eqref{eq:1.7.8}?

\item A more realistic assumption about the rental price of capital may be that there is
  an economy-wide capital market so $p_{i2}$ is the same across firms.  In this case,
  \begin{enumerate}[label=(\alph*)]
  \item Can we estimate the technology parameters?  \textbf{Hint:} The answer is yes, but
    why?  When $p_{i2}$ is constant,~\eqref{eq:1.7.4} will have the perfect
    multicollinearity problem.  But recall that $(\beta_1, \ldots, \beta_5)$ are not free
    parameters.
  \item Can we test homogeneity of the cost function in factor prices?
  \end{enumerate}

\item Taking logs of both sides of the production function~\eqref{eq:1.7.1}, one can
  derive the log-linear relationship:
  \[
    \log(Q_i) = \alpha_0 + \alpha_1\log(x_{i1}) + \alpha_2\log(x_{i2})
      + \alpha_3\log(x_{i3}) + \varepsilon_i,
  \]
  where $\alpha_0 \equiv \E[\log(A_i)]$ and
  $\varepsilon_i \equiv \log(A_i) - \alpha_0$.
\end{enumerate}
